\section{Method}

As shown in Figure~\ref{fig:pipeline}, our pipeline sequentially enriches the track catalog with external metadata, routes dialogue through dual text-matching and conversational retrieval pools, calibrates candidate rankings via LightGBM, and structures the final output using an evidence-grounded PAS framework.

\subsection{Official Data and External Metadata}
% 데이터 풍부화 및 메타데이터 결합
% 간단 설명: 시스템이 사용하는 기본 데이터셋을 명시하고, 검색과 답변 생성의 품질을 높이기 위해 외부 메타데이터를 어떻게 결합하여 데이터를 풍부하게(Enrichment) 만들었는지 설명하는 섹션입니다.
% - 기본 데이터: 공식 Challenge 데이터셋의 규모(47,071개 트랙)와 포함된 기본 정보(대화 메시지, 대화 목적 등).
% - 외부 소스 결합: 데이터 공백을 메우기 위한 3가지 외부 소스 명시 (TalkPlayTools-Env, LRCLIB, MusicBrainz) 및 각 소스의 역할.
% - 통계 제시 (Table 1 관련): 각 메타데이터의 커버리지 비율 제시 (적어도 하나 이상의 외부 메타데이터를 확보한 비율이 96.6%에 달함을 강조).
% - 구체적 예시 파싱: "Dollhouse" 사례를 통해 기본 태그에 외부 메타데이터(BPM, Key, 레이블 등)가 어떻게 덧붙여지는지 설명.
% - 피처 할당 전략: 확장된 필드들이 뒤쪽 모듈에서 어떻게 분기되어 쓰이는지 명시 (Keywords ➔ BM25, 음악 속성 ➔ 대화용 brief, 텍스트 증거 ➔ PAS 생성기).
% - 결측치 처리: 없는 데이터는 인위적으로 채워 넣지(Imputation) 않고 제외하는 미니멀한 처리 방식 언급.

The official TalkPlayData-Challenge dataset suite\footnote{\url{https://huggingface.co/datasets/talkpl-ai/TalkPlayData-Challenge}} supplies 47,071 track records with dialogue messages and conversation goals. To mitigate data sparsity, we integrate external metadata from TalkPlayTools\allowbreak-Env\footnote{\url{https://huggingface.co/datasets/talkpl-ai/TalkPlayTools-Env}} for descriptive audio, LRCLIB\footnote{\url{https://lrclib.net/}} for lyrics fallback, and MusicBrainz\footnote{\url{https://musicbrainz.org/}} for structured fields. As shown in Table~\ref{tab:metadata}, these sources provide complementary coverage, enriching 96.6\% of the catalog with at least one external field.

\begin{table}[h]
  \caption{Coverage of external metadata over the 47,071-track catalog.
  Source rows may overlap.}
  \label{tab:metadata}
  \centering
  \scriptsize
  \setlength{\tabcolsep}{2.5pt}
  \begin{tabular}{@{}lp{2.85cm}rr@{}}
    \toprule
    Source & Information used & Tracks & Coverage \\
    \midrule
    TalkPlayTools-Env & lyrics, caption, chord, tempo, key & 13,745 & 29.2\% \\
    LRCLIB & lyrics fallback & 28,162 & 59.8\% \\
    MusicBrainz & IDs, release, label, country, tags & 38,585 & 82.0\% \\
    Any source & at least one sourced field & 45,475 & 96.6\% \\
    \bottomrule
  \end{tabular}
\end{table}

To illustrate this enrichment, the track ``Dollhouse'' shown in Figure~2b expands its baseline gothic and industrial tags with external fields like 130.21 BPM and D\# minor. This enriched profile routes directly downstream: keywords expand BM25 documents, structured properties populate conversational track briefs, and textual captions and lyrics ground PAS response generation, strictly omitting missing fields.

\subsection{Pool-Specific Retrieval Interfaces}
% Hybrid & Conversational 듀얼 풀 검색

% Example provenance (not rendered): public test split
% session_id=683d0163-d246-4b08-8184-b35976816d96, turn=6.
Rather than a single query, our architecture routes retrieval through two parallel interfaces (Figure~\ref{fig:textviews}). The hybrid pool (\texttt{bm25\_qmr}) combines sparse and dense mechanisms. The sparse track runs two BM25 searches~\cite{robertson1994okapi}: a baseline catalog search and a variant enriched with catalog tags and external keywords. Its query fuses prior music turn metadata with filtered current-request tokens, and an RRF layer~\cite{cormack2009rrf} merges their rankings. Concurrently, the dense track uses a 0.6B Qwen3 embedding model~\cite{qwen3embedding2025}, joining historical metadata with the lowercased current request to query official document vectors. A second RRF fuses the sparse and dense rankings without score calibration, producing the final candidate set.


\begin{figure*}[t]
  \centering
  {\sffamily\scriptsize\bfseries (a) Final pool topology}\par\vspace{1pt}
  \begin{tikzpicture}[node distance=3.5mm and 4.5mm]
    \node[figinput,align=center,text width=18mm,inner sep=2.5pt] (hybridin)
      {Dialogue history\\+ track catalog};
    \node[figinterface,align=left,text width=59mm,inner sep=2.5pt,right=of hybridin] (hybridlogic)
      {\textbf{BM25:} catalog + external keywords $\rightarrow$ two searches $\rightarrow$ RRF\\
       \textbf{0.6B dense:} catalog strings $\rightarrow$ embedding search};
    \node[figprocess,align=center,text width=19mm,inner sep=2.5pt,right=of hybridlogic] (hybridrrf)
      {cross-component\\RRF};
    \node[figoutput,align=center,text width=27mm,inner sep=2.5pt,right=of hybridrrf] (hybridout)
      {\textbf{Hybrid pool}\\\texttt{bm25\_qmr}};

    \node[figinput,align=center,text width=18mm,inner sep=2.5pt,below=of hybridin] (convin)
      {Dialogue history\\+ track catalog};
    \node[figinterface,align=center,text width=34mm,inner sep=2.5pt,right=of convin] (convlogic)
      {8B structured query\\+ track brief};
    \node[figprocess,align=center,text width=27mm,inner sep=2.5pt,right=of convlogic] (foldsearch)
      {five adapted 8B searches};
    \node[figprocess,align=center,text width=16mm,inner sep=2.5pt,right=of foldsearch] (foldrrf)
      {five-fold\\RRF};
    \node[figoutput,align=center,text width=31mm,inner sep=2.5pt,right=of foldrrf] (convout)
      {\textbf{Conversational pool}\\\texttt{qemb\_twotower\_8b}};

    \draw[-{Latex[length=2mm]},draw=FigNavy,line width=0.9pt] (hybridin.east) -- (hybridlogic.west);
    \draw[-{Latex[length=2mm]},draw=FigNavy,line width=0.9pt] (hybridlogic.east) -- (hybridrrf.west);
    \draw[-{Latex[length=2mm]},draw=FigNavy,line width=0.9pt] (hybridrrf.east) -- (hybridout.west);
    \draw[-{Latex[length=2mm]},draw=FigNavy,line width=0.9pt] (convin.east) -- (convlogic.west);
    \draw[-{Latex[length=2mm]},draw=FigNavy,line width=0.9pt] (convlogic.east) -- (foldsearch.west);
    \draw[-{Latex[length=2mm]},draw=FigNavy,line width=0.9pt] (foldsearch.east) -- (foldrrf.west);
    \draw[-{Latex[length=2mm]},draw=FigNavy,line width=0.9pt] (foldrrf.east) -- (convout.west);
  \end{tikzpicture}

  \vspace{6pt}
  {\sffamily\scriptsize\bfseries (b) Shortened turn-6 query/document views with temporal alignment}\par\vspace{2pt}
  \scriptsize
  \setlength{\tabcolsep}{0pt}
  \renewcommand{\arraystretch}{1.3} % [핵심] 표 안의 줄 간격 및 위아래 선과의 여백을 30% 넓혀줌
  
  % 1. 메인 컴포넌트 테이블
  \begin{tabular*}{\linewidth}{@{\extracolsep{\fill}}lp{0.36\linewidth}p{0.34\linewidth}p{0.14\linewidth}@{}}
    \toprule
    \textbf{Component} & \textbf{Query-side text (Inputs)} & \textbf{Document-side text (Target)} & \textbf{Preserved signal} \\
    \midrule
    \textbf{Hybrid: BM25} &
    \textbf{[Turns 1--5]} \texttt{Burn Witch Burn; Ego Likeness; 2006} \newline 
    \textbf{[Turn 6]} \texttt{gothic rock; female vocals; industrial; driving beat} &
    \texttt{Dollhouse (Razed In Black Mix); Switchblade Symphony; 2001}; tags: \texttt{gothic, darkwave, industrial}. & 
    Exact entities and current search tags. \\
    \midrule % hline을 midrule로 교체하여 숨통을 틔움
    \textbf{Hybrid: 0.6B dense} &
    \textbf{[Turns 1--5]} \texttt{title: burn witch burn, artist: ego likeness} \newline 
    \textbf{[Turn 6]} Lowercase text-based search tokens. & 
    \texttt{Dollhouse; Switchblade Symphony; 2001; popularity 11}; tags: \texttt{punk, ebm, remix}. & 
    Catalog semantics; official fields only. \\
    \midrule
    \textbf{Conversational: 8B} &
    Fixed four-field structured prompt template. \newline (Instantiated in the table below.) & 
    \texttt{Dollhouse} by Switchblade Symphony; 130 BPM, D\# minor, gothic/darkwave. & 
    Evolving conversational intent across turns. \\
    \bottomrule
  \end{tabular*}
  
  % 첫 번째 테이블 오른쪽 아래 풋노트 배치
  % \vspace{1pt}
  % \hfill \textit{$[\dagger]$ Instantiated in the table below.}\hspace{2pt}

  % \vspace{8pt}

  % 2. 하단 8B 상세 구조화 쿼리 테이블 (1단 압축 레이아웃)
  \begin{tabular*}{\linewidth}{@{\extracolsep{\fill}} p{0.22\linewidth} p{0.74\linewidth} @{}}
    \toprule
    \multicolumn{2}{@{}l}{\textbf{Actual 8B Structured Query Fields for Turn-6 Execution}} \\
    \midrule
    \textbf{History [Turns 1--5]} & 
    \textbf{Recent listens:} Burn Witch Burn [T5]; In The Throat Of The Unbounded; Sorrow Is Her Name; Quiet Moments; Frozen. \newline
    \textbf{Session so far:} User anchored in Lycia's ethereal darkwave, expanding through Violet Tears and Ego Likeness toward gothic rock. \\
    \midrule
    \textbf{Current [Turn 6]} & 
    \textbf{Looking for:} gothic rock + industrial elements + heavier driving beat. \newline
    \textbf{Context:} female vocals; Switchblade Symphony; gothic/darkwave. \\
    \bottomrule
  \end{tabular*}
  \vspace{-3mm}
  \caption{Two final-pool pipelines and temporal separation of inputs. (a) Topology of the hybrid and conversational pipelines. (b) Input-target pairs for Turn 6. The bottom subtable details the 8B model's query structure.}
  \label{fig:textviews}
  \vspace{-2mm}
\end{figure*}

The conversational pool (\texttt{qemb\_twotower\_8b}) relies entirely on natural language on both sides to track long-range session dynamics. The query format strictly structures the dialogue history into four ordered fields: \texttt{Recent listens:} lists prior tracks chronologically; \texttt{Session so far:} supplies an optional cached preference trajectory summary under 25 words; \texttt{Context:} aggregates explicit goals, eras, and popularity cues; and \texttt{Looking for:} contains the current explicit request. When context window limits are reached, a strict truncation priority rule preserves the user's immediate intent: because \texttt{Looking for:} sits adjacent to the final embedding token, it is never dropped, whereas lower-priority historical context fields are truncated first. To optimize this contextual mapping, we fine-tune an 8B parameter encoder model using 4-bit QLoRA~\cite{dettmers2023qlora} under an in-batch InfoNCE contrastive objective~\cite{oord2018cpc} with LogQ frequency correction~\cite{yi2019sampling}. The dense retrieval is executed across five folds, each optimized for distinct projections of track briefs and user embedding spaces. Their rankings are then combined via RRF into the unified \texttt{qemb\_twotower\_8b} pool.


Turn 1 is handled as an exception due to the lack of dialogue history. In the hybrid pool, both BM25 and dense queries collapse into the current request tokens. In the conversational pool, the 8B query omits \texttt{Recent listens:} and \texttt{Session so far:} but retains goal cues and \texttt{Looking for:}. While document representations remain unchanged, the downstream LightGBM feature row records the turn number. This enables the subsequent ranking and calibration stages to distinguish cold-start turns from multi-turn dialogue contexts.


\subsection{Candidate Fusion and Calibration}
% LightGBM 기반 순위 보정 및 융합

To consolidate retrieval, the system builds a top-100 union from both pools. A LightGBM~\cite{ke2017lightgbm} ranker then scores and calibrates this final list. Table~\ref{tab:features} outlines the input feature schema. The model ingests individual ranks and scores from both pools, alongside their cross-pool RRF metrics. Turn number, request category, and goal specificity define the session context. Non-retrieved candidates receive sentinel values. Training utilizes out-of-fold scoring to eliminate data leakage from the 8B retrieval model.

\begin{table}[h]
  \caption{Released LightGBM feature schema.}
  \label{tab:features}
  \centering
  \scriptsize
  \setlength{\tabcolsep}{3pt}
  \begin{tabular}{@{}lll@{}}
    \toprule
    Source & Fields & Decision evidence \\
    \midrule
    Hybrid pool & rank, score & exact/general support \\
    Conversational pool & rank, score & dialogue support \\
    Cross-pool RRF & rank, score & retrieval agreement \\
    Context & turn, category, specificity & request regime \\
    \bottomrule
  \end{tabular}
\end{table}

\begin{figure*}[t]
  \centering
  \begin{tikzpicture}[node distance=3mm,figure-readable]
    \node[figinput,align=left,text width=28mm,inner sep=3pt,font=\sffamily\tiny] (ev)
      {\textbf{Evidence}\par
       Request + dialogue\par
       Top-5 details + closed top-20 title list\par
       \vspace{1pt}\textcolor{FigLine}{\rule{\linewidth}{0.4pt}}\par\vspace{1pt}
       \textcolor{FigTab}{\bfseries Request evidence}\par
       ``driving bassline''; ``dark and intense''; prior Flying Lotus preference};
    \node[figprocess,align=left,text width=28mm,inner sep=3pt,font=\sffamily\tiny,
      right=2.5mm of ev.north east,anchor=north west] (prop)
      {\textbf{Propose}\par
       Assign each candidate to its first qualifying group under fixed precedence\par
       \vspace{1pt}\textcolor{FigLine}{\rule{\linewidth}{0.4pt}}\par\vspace{1pt}
       \textcolor{FigTab}{\bfseries Candidate-level groups}\par
       ``Camel'' similarity: 0.60 $\geq$ threshold\par
       $\rightarrow$ query-aligned};
    \node[figfocus,align=left,text width=43mm,inner sep=3pt,font=\sffamily\tiny,
      right=2.5mm of prop.north east,anchor=north west] (llm)
      {\textbf{LLM generation}\hfill\textcolor{FigTab}{\bfseries ONE CALL}\par
       Use proposed groups + one routed demo to perform Assign + Select and emit
       the response\par
       \vspace{1pt}\textcolor{FigLine}{\rule{\linewidth}{0.4pt}}\par\vspace{1pt}
       \textcolor{FigTab}{\bfseries Generation fields}\par
       core = ``Camel''; reference = prior Flying Lotus; supplement = none;
       theme = ``darker, rawer pocket''};
    \node[figprocess,align=left,text width=31mm,inner sep=3pt,font=\sffamily\tiny,
      right=2.5mm of llm.north east,anchor=north west] (sel)
      {\textbf{Validate / repair}\par
       Check titles, claims, supplements, and repeated phrasing\par
       \vspace{1pt}\textcolor{FigLine}{\rule{\linewidth}{0.4pt}}\par\vspace{1pt}
       \textcolor{FigTab}{\bfseries Validation result}\par
       title allowed; claim supported; no supplement; soft-pattern checks pass};
    \node[figoutput,align=left,text width=24mm,inner sep=3pt,font=\sffamily\tiny,
      right=2.5mm of sel.north east,anchor=north west] (resp)
      {\textbf{Grounded response}\par
       Return the verified explanation\par
       \vspace{1pt}\textcolor{FigLine}{\rule{\linewidth}{0.4pt}}\par\vspace{1pt}
       \textcolor{FigTab}{\bfseries Final excerpt}\par
       ``Camel'' sits in a darker, rawer pocket: vibrating sub-bass and
       stripped-down, menacing production.};
    \draw[figarrow] ([yshift=-9mm]ev.north east)--([yshift=-9mm]prop.north west);
    \draw[figarrow] ([yshift=-9mm]prop.north east)--([yshift=-9mm]llm.north west);
    \draw[figarrow] ([yshift=-9mm]llm.north east)--([yshift=-9mm]sel.north west);
    \draw[figarrow] ([yshift=-9mm]sel.north east)--([yshift=-9mm]resp.north west);

    \node[draw=FigTab,dashed,line width=0.7pt,fill=FigLavender,rounded corners=2pt,
      align=center,text width=52mm,inner sep=3pt,font=\sffamily\tiny,
      above=3mm of llm] (demo)
      {\textcolor{FigTab}{\bfseries Few-shot demo bank $\cdot$ Fig.~\ref{fig:demo-compiler}}\par
       retrieve a bucket-matched example within a single LLM call};
    \draw[figdash] (demo.south)--(llm.north);
  \end{tikzpicture}
  \caption{PAS runtime execution mapped to a single-shot demonstration workflow. Candidate-level evidence groups and a bucket-matched example condition a single LLM call, followed by deterministic validation.}
  \label{fig:pas}  
\end{figure*}

\begin{figure*}[t]
  \centering
  \begin{tikzpicture}[node distance=3mm,figure-readable]
    \node[figinput,align=left,text width=22mm,inner sep=2.5pt,font=\sffamily\tiny] (turns)
      {\textbf{Eligible training turns}\par
       Sample completed recommendation turns\par
       \vspace{1pt}\textcolor{FigLine}{\rule{\linewidth}{0.4pt}}\par\vspace{1pt}
       \textcolor{FigTab}{\bfseries Compiler input}\par
       1,000 turns, not sessions};
    \node[figinterface,align=left,text width=26mm,inner sep=2.5pt,font=\sffamily\tiny,
      right=2.5mm of turns.north east,anchor=north west] (judge)
      {\textbf{Generate + judge}\par
       GLM-5.2 candidate generation\par
       Personalization, Explanation, and lexical novelty\par
       \vspace{1pt}\textcolor{FigLine}{\rule{\linewidth}{0.4pt}}\par\vspace{1pt}
       \textcolor{FigTab}{\bfseries Quality gates}\par
       15--130 words; total score $\geq$ 0.7};
    \node[figprocess,align=left,text width=24mm,inner sep=2.5pt,font=\sffamily\tiny,
      right=2.5mm of judge.north east,anchor=north west] (coverage)
      {\textbf{Coverage selection}\par
       Balance category, specificity, and turn position\par
       \vspace{1pt}\textcolor{FigLine}{\rule{\linewidth}{0.4pt}}\par\vspace{1pt}
       \textcolor{FigTab}{\bfseries Selected bank}\par
       18 demonstrations};
    \node[figfocus,align=left,text width=32mm,inner sep=2.5pt,font=\sffamily\tiny,
      right=2.5mm of coverage.north east,anchor=north west] (bank)
      {\textbf{Bucketed demo bank}\par
       Index by category $\times$ specificity $\times$ turn\par
       \vspace{1pt}\textcolor{FigLine}{\rule{\linewidth}{0.4pt}}\par\vspace{1pt}
       \textcolor{FigTab}{\bfseries Matched bucket}\par
       key by A | HL | mid\par
       \textcolor{FigMuted}{compiled once; reused at runtime}};
    \node[figprocess,align=left,text width=25mm,inner sep=2.5pt,font=\sffamily\tiny,
      right=2.5mm of bank.north east,anchor=north west] (retrieve)
      {\textbf{Retrieve one demo}\par
       Select one compiled response example from the matched bucket\par
       \vspace{1pt}\textcolor{FigLine}{\rule{\linewidth}{0.4pt}}\par\vspace{1pt}
       \textcolor{FigTab}{\bfseries Selected example}\par
       core track: ``Camel'' by Flying Lotus};
    \node[figoutput,draw=FigTab,line width=1.0pt,align=left,text width=22mm,inner sep=2.5pt,
      font=\sffamily\tiny,right=2.5mm of retrieve.north east,anchor=north west] (handoff)
      {\textbf{Condition Fig.~\ref{fig:pas}}\par
       Insert bucket-matched example as one-shot response guidance
       \vspace{1pt}\textcolor{FigLine}{\rule{\linewidth}{0.4pt}}\par\vspace{1pt}
       \textcolor{FigTab}{\bfseries Destination}\par
       single LLM generation call};

    \draw[figarrow] ([yshift=-9mm]turns.north east)--([yshift=-9mm]judge.north west);
    \draw[figarrow] ([yshift=-9mm]judge.north east)--([yshift=-9mm]coverage.north west);
    \draw[figarrow] ([yshift=-9mm]coverage.north east)--([yshift=-9mm]bank.north west);
    \draw[figarrow] ([yshift=-9mm]bank.north east)--([yshift=-9mm]retrieve.north west);
    \draw[figarrow] ([yshift=-9mm]retrieve.north east)--([yshift=-9mm]handoff.north west);

    \draw[draw=FigLine,line width=0.7pt]
      ([yshift=2mm]turns.north west)--([yshift=2mm]bank.north east);
    \node[font=\sffamily\tiny\bfseries,text=FigMuted,fill=white,inner sep=1pt,
      above=2mm of $(turns.north)!0.5!(bank.north)$] {COMPILE ONCE OFFLINE};
    \draw[draw=FigTab,line width=0.7pt]
      ([yshift=2mm]retrieve.north west)--([yshift=2mm]handoff.north east);
    \node[font=\sffamily\tiny\bfseries,text=FigTab,fill=white,inner sep=1pt,
      above=2mm of $(retrieve.north)!0.5!(handoff.north)$] {RETRIEVE AT RUNTIME};

    \node[draw=FigTab,dashed,line width=0.7pt,fill=FigLavender,rounded corners=2pt,
      align=center,text width=32mm,inner sep=2.5pt,font=\sffamily\tiny,
      below=3mm of bank] (request)
      {\textcolor{FigTab}{\bfseries Current request signature}\par category =  \texttt{A} | specificity = \texttt{HL} | turn = \texttt{mid}};
    \draw[figdash] (request.north)--(bank.south);
  \end{tikzpicture}
  \caption{Offline demonstration compilation and runtime routing mechanisms. The current request signature queries the shared 18-example bank. The released \texttt{A|HL|mid} route then selects a matching example.}
  \label{fig:demo-compiler}
\end{figure*}

This framework directly resolves cross-pool disagreements. In the gothic/industrial turn-6 case, the target track ranks 20th in the hybrid pool and 4th in the conversational pool. Cross-pool RRF elevates it to 2nd place, and LightGBM reranks it to 1st. Table~\ref{tab:features} proves the feature footprint is lightweight. Semantic text reinterpretation remains entirely upstream, leaving LightGBM to act strictly as an efficient coordinator of retrieval agreement and context.

\subsection{PAS Response Generation}
\label{sec:pas_generation}
% PAS 프레임워크 기반 런타임 데모 라우팅 및 답변 생성
% 기존 원문은 런타임 단계와 오프라인 빌드 단계가 혼재되어 있고, 복잡한 리스트가 한 문장에 얽혀 있어 인지 부하가 높았습니다. 이를 데이터가 생성되고 소비되는 순서에 맞추어 다음과 같은 4단계 파이프라인 흐름으로 재구조화했습니다.

% 1. **[개념 정의 및 Propose 단계 (온라인 1단계)]**: PAS 프레임워크의 목적을 명시하고, Top-20 후보군을 고정된 우선순위 규칙에 따라 개별 증거 그룹으로 분류하는 단계를 다룹니다.
% 2. **[Assign, Select 및 사후 검증 단계 (온라인 2, 3단계)]**: 분류된 그룹 정보를 단일 LLM Call로 처리하여 최종 답변을 생성하고, 규칙 기반으로 오류를 실시간 수선하는 과정을 설명합니다.
% 3. **[퓨샷 데모 뱅크 구축 (오프라인 파이프라인)]**: 대화 맥락에 최적화된 18개의 고품질 가이드라인 샘플을 오프라인에서 엄격한 필터링 기법으로 컴파일하는 과정을 다룹니다.
% 4. **[런타임 라우팅 메커니즘 (온오프라인 연결)]**: 실시간 입력 쿼리의 특징(시그니처)을 기반으로 백업 레이어를 순차 탐색하여 최적의 데모를 매칭하는 규칙을 정립합니다.

We adapt the PAS framework~\cite{wang2023goalex} to enforce structural constraints on response claims for the top-20 ranked tracks. Figure~\ref{fig:pas} illustrates the complete runtime execution stages alongside a sample execution trace. First, the propose stage partitions these 20 candidates into distinct evidence groups based on fixed precedence rules. These groups consist of query alignment, lyric resonance, taste continuity, pool coherence, and discovery. In the sample trace, the text tokens corresponding to driving bassline and dark and intense produce a 0.60 similarity score for the track Camel by Flying Lotus. This score successfully satisfies the established query alignment threshold.

A single structured LLM call processes the proposed candidate groups together with one routed demonstration. This demonstration provides response selection guidance rather than introducing additional recommendation candidates. Second, within this single call, the assign step identifies usable personalization anchors and supported themes. Third, the select step then chooses a core track, optional supplementary tracks for plural requests, and optional historical references. The same joint call generates the raw response text. Lastly, deterministic rules validate and repair title formatting errors or citation violations. If the draft matches any blacklisted session phrases, PAS triggers a single generation retry.

Figure~\ref{fig:demo-compiler} outlines the offline compilation process for the reusable demonstration bank. The compilation pipeline samples 1,000 eligible recommendation turns from the TalkPlayData training split. GLM-5.2 generates the initial candidate responses. An independent LLM-as-judge scores personalization and explanation metrics on a 1-to-5 scale. We normalize these outputs into a unified quality metric. The final evaluation combines 75\% of this quality metric with 25\% lexical novelty. Viable candidates must contain between 15 and 130 words and achieve a minimum combined score of 0.7. A coverage-aware filtering step balances the final selection across request category, specificity, and turn position. This selection process yields a fixed 18-example runtime demonstration bank.

At runtime, each incoming request generates a unique routing key based on its category, specificity, and turn. The router evaluates the exact matching bucket first. If that bucket is empty, the system sequentially falls back to category matches, specificity or turn matches, and the global bank. The router injects the highest-scoring demonstration from the first nonempty layer. Figure~\ref{fig:demo-compiler} displays the key \texttt{A|HL|mid}, where \texttt{A} represents a sound-focused category, \texttt{HL} denotes the specificity regime, and \texttt{mid} indicates a mid-dialogue turn. The concrete trace details mapped in Figure~\ref{fig:pas} reflect this specific routed instance to ground the respective PAS input fields.